\documentclass{beamer}
%% fink install texlive
\usetheme{Warsaw}
\usepackage{amsmath,amssymb}
\usepackage{stmaryrd}
\usepackage{graphicx}
%\usepackage{haskell}
\usepackage{code}
%\usepackage{proof}
\usepackage{theorem}
\usepackage{pstricks} 


\usepackage{tikz}
\usetikzlibrary{shapes.geometric, arrows, shapes}

\tikzstyle{startstop} = [rectangle, rounded corners, minimum width=1.5cm, minimum height=0.5cm,text centered, draw=black]
\tikzstyle{io} = [trapezium, trapezium left angle=70, trapezium right angle=110, minimum width=1.5cm, minimum height=0.5cm, text centered, draw=black]
\tikzstyle{process} = [rectangle, minimum width=1.5cm, minimum height=0.5cm, text centered, draw=black]
\tikzstyle{decision} = [diamond, minimum width=1.5cm, minimum height=0.5cm, text centered, draw=black, aspect=3]
\tikzstyle{arrow} = [thick,->,>=stealth]
\tikzstyle{doublearrow} = [thick,<->,>=stealth]
\tikzstyle{db} = [cylinder, shape border rotate=90, draw,minimum
height=1.5cm, minimum width=2cm, aspect=0.3]
\tikzstyle{www} = [cloud, cloud puffs=15.7, cloud ignores
aspect,minimum width=5cm, minimum height=2cm, align=center, draw]
\tikzstyle{blank} = [rectangle] 
\tikzstyle{container} = [rectangle, rounded corners, minimum
width=7cm, minimum height=7cm,text centered, draw=black]

\tikzstyle{darrow} = [thick,<->,>=stealth]

\usepackage{listings}

\include{macros-ms}

\newcommand{\rev}[1]{{#1}^{\mbox{\scriptsize r}}}
\newcommand{\lquo}{\backslash}
\newcommand{\rquo}{/}
\newcommand{\diff}{-}
\newcommand{\isect}{\cap}
\newcommand{\mymid}{~\redtxt{\mid}~}
\newcommand{\lconf}[2]{#1:{\tt #2}}

\newcommand{\ignore}[1]{}

\newcommand{\magtxt}[1]{{\magenta #1}}
\newcommand{\redtxt}[1]{{\red #1}}
\newcommand{\bluetxt}[1]{{\blue #1}}
\newcommand{\greytxt}[1]{{\gray #1}}
\newcommand{\greentxt}[1]{{\green #1}}
\newcommand{\mmleq}{\leq}
\newcommand{\pow}{\^{}}
\newcommand{\venv}{\Delta}
\newcommand{\mleq}{\mbox{\tt leq}}
\newcommand{\mas}{\mbox{\tt as}}
\newcommand{\subt}{<\!:}

\newcommand{\Nturns}{\, \vdash_{\mbox{\tiny lnf}} \,}

\newcommand{\deriv}[2]{#1\backslash #2}
\newcommand{\pderiv}[2]{pd_{#2}(#1)} %% {#1\backslash_p #2}
\newcommand{\pderivFail}[2]{pd^{err}_{#2}(#1)} %% {#1\backslash_p #2}
\newcommand{\mkEmp}[1]{{#1}_{\epsilon}}
\newcommand{\strip}[1]{#1\downarrow}

\newenvironment{ttprog}{\begin{trivlist}\item \tt
        \begin{tabbing}}{\end{tabbing}\end{trivlist}}


\newenvironment{grammar}{%
         \begin{center} \small%
         $\begin{array}{rcll}
         }{%
         \end{array}$\end{center}\ignorespaces%
         }

\newcommand{\rem}[1]{}
\newcommand{\comment}[1]{}

\newcommand{\kl}[1]{}

\bibliographystyle{plainnat}

\begin{document}
%\lstset{language=python}
\logo{
\includegraphics[height=0.8cm]{pics/nyp-logo.eps}\vspace{230pt}
\includegraphics[height=0.6cm]{pics/sit-logo.png}\vspace{230pt}
}
\title{Introduction to Scala} 
\author{
 Kenny Zhuo Ming Lu
}
\institute[Inst.]{School of Information Technology, Nanyang Polytechnic}
\date{\today} 


\frame{\titlepage} 

%\bibliographystyle{plainnat}
%-------------------------------------------------------------------
%-------------------------------------------------------------------
\begin{frame}[fragile]
\frametitle{Agenda}
\begin{itemize}
 \item Introduction
 \item Object Oriented Programming 
 \item Functional Programming 
 \item (Generalized) Algebraic Data Type and pattern matching
 \item Functional Abstraction and Higher order function
 \item Type classes and Monads
 \item 
\end{itemize}
\end{frame}

%-------------------------------------------------------------------
%-------------------------------------------------------------------
\begin{frame}[fragile]
\frametitle{Scala is ... }

A general purpose programming language that is 
\begin{itemize}
 \item Object Oriented and Functional
 \item with distributed programming builtin
 \item with a rich type system and a huge set of libraries and eco-systems
 \item Widely used in research institutes (e.g. Spark) and industries
   (e.g. Twitter, Linkedin, EDF, ... )
\end{itemize}
\end{frame}



\end{document}

